
% Add "answers" option to print the answers
\documentclass[addpoints,answers]{exam}

\usepackage[utf8]{inputenc}
\usepackage{calc}
\usepackage{listings}
\usepackage{xcolor}
\usepackage{xspace}

\definecolor{codegreen}{rgb}{0,0.6,0}
\definecolor{codegray}{rgb}{0.5,0.5,0.5}
\definecolor{codepurple}{rgb}{0.58,0,0.82}

\lstdefinestyle{Cstyle}{
    commentstyle=\color{codegreen},
    keywordstyle=\color{magenta},
    numberstyle=\tiny\color{codegray},
    stringstyle=\color{codepurple},
    basicstyle=\ttfamily\small,
    breakatwhitespace=true,         
    breaklines=true,                 
    captionpos=b,                    
    keepspaces=true,                 
    numbers=left,                    
    numbersep=5pt,                  
    showspaces=false,                
    showstringspaces=false,
    showtabs=false,                  
    tabsize=4
}

\lstset{style=Cstyle}

% \makebox[\textwidth]{Name:\enspace\hrulefill}
%  
% \vspace{5mm}
%\makebox[\textwidth]{Instructor’s name:\enspace\hrulefill}

\newcommand{\truefalse}
{		\begin{checkboxes}
			 \choice True
			 \choice False. Explanation: 
		\end{checkboxes}
		\vspace{10pt}
}

\newcommand{\answerspace}[1]
{
	\ifprintanswers
	\else
		\setlength{\mylength}{10pt}%
		%\noindent\rule{\mylength}{20pt}
		%\vspace{\fpeval{1.2*(#1)}}
		\vspace{#1\mylength}
	\fi
}
\newlength{\mylength}

\newcommand\pf[1]{\textcolor{red}{PF: #1}}
\newcommand\am[1]{\textcolor{blue}{AM: #1}}

\newcommand{\answer}[1]{\begin{solution}{#1}\end{solution}}
\newcommand{\ra}{$\rightarrow$\xspace}



\usepackage{color}
\usepackage{lastpage}
\usepackage{verbatim}
\usepackage{fancyvrb}
\usepackage{graphicx}
\usepackage{float}
\usepackage{algorithm}
\usepackage{listings}
\usepackage{algpseudocode}
\usepackage{tabularx}
\usepackage{ulem}

\newcommand{\includecode}[2][c]{\lstinputlisting[caption=, escapechar=, language=#1]{#2}<!---->}
\RecustomVerbatimCommand{\VerbatimInput}{VerbatimInput}{}

\renewcommand{\solutiontitle}{\noindent\textbf{Solution sketch: }}

\definecolor{SolutionColor}{rgb}{0.8,0.9,1}


\cfoot{Page \thepage\ of \pageref{LastPage}}

\rhead{Initials: \rule{20mm}{0.4pt}}



%%%%%%%%%%%%%%%%%%%%%%%%%%%%%%%%%%%%%%%%
\begin{document}
\fvset{numbers=left,numbersep=3pt}
\thispagestyle{empty}

\begin{center}
	{\LARGE CS408 Fall 2025 -- Homework 5 (77 points)\\}

	{\Large Deadline: December 5, Friday\\}

	\vspace{20mm}
\fbox{\fbox{\parbox{5.5in}{\centering
	\begin{itemize}
		\item The homework is open notes, open book.
		\item No collaboration and no communication of any kind allowed.
		\item Provide complete and clear but concise answers. Justify thoroughly statements where required.
		\item Coding questions should be answered in C.
		\item Use basic blocks for all CFG exercises.
		\item For questions related to compilation, test code on the Purdue ``data'' server and with the default gcc compiler on the ``data'' server.
	\end{itemize}
	}}}
\end{center}

\vspace{13mm}
\begin{center}
	\makebox[0.6\textwidth]{
		\combinedgradetable[h][questions]
	}
\end{center}
%\begin{center}\combinedgradetable[h][questions]\end{center}
\vspace{28mm}

\begin{center}
\makebox[\textwidth]{\textbf{Name (please print)}: \rule{100mm}{0.8pt} }
\end{center}

\newpage
\begin{questions}

%-----Question 1-----%
\question Answer the following questions related to Integration testing:
\begin{parts}
	\part[4] Please provide the advantages of Bottom Up Integration testing over Top Down Integration testing.
	\answer{
		\answerspace{10}
		\begin{itemize}
			\item \textbf{Early testing of low-level modules}: Bottom Up Integration testing allows for the early testing of low-level modules, which can help identify issues in foundational components before integrating them into higher-level modules.
			\item \textbf{No need for stubs}: In Bottom Up Integration testing, there is no need to create stubs for higher-level modules, which can simplify the testing process and reduce the amount of test code that needs to be written.
		\end{itemize}
	}
	\vspace{5mm}

	\part[4] Please provide the advantages of Top Down Integration testing over Bottom Up Integration testing.
	\answer{
		\answerspace{10}
		\begin{itemize}
			\item \textbf{Early detection of high-level design flaws}: Top Down Integration testing allows for the early detection of design flaws in high-level modules, which can be critical for the overall architecture of the system.
			\item \textbf{User interface testing}: It enables testing of the user interface early in the development process, ensuring that the user experience is considered from the start.
		\end{itemize}
	}
\end{parts}
\vspace{5mm}

%-----Question 2-----%
\question Consider a hypothetical \textbf{Library Management System} designed with the following three-layer architecture:

\begin{itemize}
    \item \textbf{Layer I (Top):} \texttt{LibraryDashboard} (User Interface)
    \item \textbf{Layer II (Middle):} \texttt{BookService} (Business Logic), \texttt{MemberService} (Business Logic)
    \item \textbf{Layer III (Bottom):} \texttt{SQLDatabase} (Data Storage), \texttt{EmailNotifier} (External Alert System)
\end{itemize}

\noindent \textbf{Dependency Rules:}
\begin{itemize}
    \item \texttt{LibraryDashboard} calls \texttt{BookService} and \texttt{MemberService}.
    \item \texttt{BookService} calls \texttt{SQLDatabase}.
    \item \texttt{MemberService} calls both \texttt{SQLDatabase} and \texttt{EmailNotifier}.
\end{itemize}

\vspace{5mm}
\begin{parts}
	\part[3] You are performing \textbf{Top-Down Integration} on this Library Management System. 
	\begin{itemize}
		\item You have successfully verified the top layer (\texttt{LibraryDashboard}) in isolation.
		\item You are now moving to the next step: \textbf{Integrating \texttt{MemberService}}.
		\item Note: \texttt{BookService} has \textit{not} been integrated yet.
	\end{itemize}
	
	To successfully execute this specific integration step without crashing, what distinct testing components (stubs or drivers) are required?
	\vspace{2mm}
	
	\begin{enumerate}
		\item[A)] 1 driver for \texttt{LibraryDashboard} and 1 stub for \texttt{BookService}.
		\item[B)] 2 stubs simulating \texttt{SQLDatabase} and \texttt{EmailNotifier}.
		\item[C)] 1 stub simulating \texttt{SQLDatabase} only.
		\item[D)] No stubs are needed because \texttt{MemberService} contains the logic.
	\end{enumerate}

	\answer{
		B
		\\
		\textbf{Explanation:} 
		\begin{enumerate}
			\item \textbf{Current context:} You are utilizing a Top-Down strategy. Layer I is present, and you are currently testing the integration of Layer II (\texttt{MemberService}).
			\item \textbf{Dependency analysis:} The \texttt{MemberService} module depends on calls to Layer III components to function (specifically \texttt{SQLDatabase} to save data and \texttt{EmailNotifier} to send alerts).
			\item \textbf{Missing components:} In a Top-Down approach, the bottom layers (Layer III) are developed or integrated \textit{last}. Therefore, the real \texttt{SQLDatabase} and \texttt{EmailNotifier} are not yet available.
			\item \textbf{The Solution:} To prevent the system from crashing when \texttt{MemberService} attempts to invoke these missing dependencies, you must create \textbf{Stubs} for both.
			\item \textbf{Why not A?} A driver for \texttt{LibraryDashboard} is unnecessary because the real module is already present and acts as the caller. A stub for \texttt{BookService} is unnecessary for this specific test because \texttt{MemberService} does not call \texttt{BookService} (they are independent siblings in Layer II).
		\end{enumerate}
	}

	\vspace{5mm}
	\part[3] You are performing \textbf{Bottom-Up Integration}. 
	\begin{itemize}
		\item You have already successfully tested and integrated all modules in Layer III (\texttt{SQLDatabase} and \texttt{EmailNotifier}).
		\item You are now ready to test \textbf{\texttt{MemberService}} (Layer II).
	\end{itemize}
	\vspace{2mm}
	Which of the following describes the correct test configuration for this step?
	
	\begin{enumerate}
		\item[A)] You need a driver simulating \texttt{LibraryDashboard} to call \texttt{MemberService}; the real Layer III modules are used.
		\item[B)] You need a stub for \texttt{SQLDatabase} and a driver for \texttt{LibraryDashboard}.
		\item[C)] You need stubs for both \texttt{SQLDatabase} and \texttt{EmailNotifier} because Layer II is the target.
		\item[D)] No stubs or drivers are needed; \texttt{MemberService} can be tested in isolation.
	\end{enumerate}
	\answer{
		A
		\\
		\textbf{Explanation:} 
		Both \texttt{BookService} and \texttt{MemberService} directly depend on (call) \texttt{SQLDatabase}. Any change in the database structure or behavior directly impacts them. \texttt{EmailNotifier} is a sibling in Layer III and does not call \texttt{SQLDatabase}, so it is unaffected.
	}


	
\end{parts}

\vspace{5mm}
%-----Question 3-----%
\question[30] Design a function in C that adheres to LOC-3 compliance rules for code written at JPL. The function should calculate the factorial of a given non-negative integer ($N!$) and return the result. \\
Please make sure to include any required C headers.\\
Ensure that the function follows these guidelines:
\begin{enumerate}
    \item It passes the compiler and a static source code analyzer without triggering warnings (Assume \texttt{gcc -Wall} satisfies both requirements).
    \item It contains a statically verifiable upper bound on the number of loop iterations (Assume that the verifier can only analyze bounds of the form "i $<$ C" in the for loop, with i in the iteration count and C a constant).
    \item The function has at least one assertion to check the input value, and the code reaches a minimal assertion density of 2\%.
\end{enumerate}
Provide a brief explanation of how your implementation follows each of these guidelines. \\
Keep your code simple and concise, up to 30 lines of code. \\
The function's signature is given below: \\
uint32\_t \_safe\_factorial(uint32\_t n);

\answerspace{10}
\answer {
	\VerbatimInput{hw5/q2.c}
	
		\textbf{Explanation of Guidelines:}
		\begin{enumerate}
			\item The code is standard C, includes \texttt{stdint.h} for types and \texttt{assert.h} for the assertion. It compiles without warnings.
			\item The loop includes a statically verifiable constant bound: $i < 13$. The verifier can easily confirm the loop will never run more than 13 times, satisfying the safety constraint regardless of the variable $n$.
			\item The assertion \texttt{assert(n <= 12)} ensures the input will not cause an integer overflow. With approximately 12 lines of code and 1 assertion, the density is $\approx 8\%$, exceeding the 2\% requirement.
		\end{enumerate}
	
}
\vspace{5mm}

%-----Question 4-----%
\question Consider three inputs, with the following representative values from each block of a partition: \\
\textbf{Partition A}: \{Block a1, Block a2\} \\
\textbf{Partition B}: \{Block b1, Block b2, Block b3\} \\
\textbf{Partition C}: \{Block c1, Block c2\, Block c3\} \\

For each of the following parts, please ensure each set entry is printed on a different line and lines are sorted.
Eg. In Lecture 32 slide 23, the answer should be

TR = \{ \\
(A, 1, x) \\
(A, 1, y) \\
(A, 2, x) \\
(A, 2, y) \\
(A, 3, x) \\
(A, 3, y) \\
(B, 1, x) \\
(B, 1, y) \\
(B, 2, x) \\
(B, 2, y) \\
(B, 3, x) \\
(B, 3, y) \\
\}

You can use the Linux utility \texttt{sort} to sort the entries.
\texttt{sort} expects one entry per line.
\vspace{5mm}
\begin{parts}
\part[6] List the TR for Each Choice Coverage (ECC).
\answer {
	TR = \\
	\{ \\
		a1 \\
		a2 \\
		b1 \\
		b2 \\
		b3 \\
		c1 \\
		c2 \\
		c3 \\
	\}
}
\vspace{5mm}
\part[9] List the TR for Pair-Wise Coverage (PWC).
\answer{
	% A B
	% A C
	% B C
	TR = \\
	\{ \\
	(a1, b1) \\
	(a1, b2) \\
	(a1, b3) \\
	(a1, c1) \\
	(a1, c2) \\
	(a1, c3) \\
	(a2, b1) \\
	(a2, b2) \\
	(a2, b3) \\
	(a2, c1) \\
	(a2, c2) \\
	(a2, c3) \\
	(b1, c1) \\
	(b1, c2) \\
	(b1, c3) \\
	(b2, c1) \\
	(b2, c2) \\
	(b2, c3) \\
	(b3, c1) \\
	(b3, c2) \\
	(b3, c3) \\
	
	\}
}

% \part[10] List the TR for Base Choice Coverage (BCC).
% \answer {
% 	For a base choice (a1, b1, c1), the TR is: \{(a2, b1, c1), (a3, b1, c1), (a1, b2, c1), (a1, b1, c2)\}
% }
\vspace{5mm}
\part[9] List the TR for All Combinations Coverage (ACoC).
\answerspace{10}
\answer {
	TR = \\
	\{ \\
	(a1, b1, c1) \\
	(a1, b1, c2) \\
	(a1, b1, c3) \\
	(a1, b2, c1) \\
	(a1, b2, c2) \\
	(a1, b2, c3) \\
	(a1, b3, c1) \\
	(a1, b3, c2) \\
	(a1, b3, c3) \\
	(a2, b1, c1) \\
	(a2, b1, c2) \\
	(a2, b1, c3) \\
	(a2, b2, c1) \\
	(a2, b2, c2) \\
	(a2, b2, c3) \\
	(a2, b3, c1) \\
	(a2, b3, c2) \\
	(a2, b3, c3) \\
	\}
}
\end{parts}

\vspace{5mm}
%-----Question 5-----%
\question[3] LOC-2 aims to secure "predictable execution" in an embedded context. When defining the upper bound for loop iterations, what specific constraint does the rule place on \textit{how} this bound must be determined?
\begin{enumerate}
    \item[\textbf{A)}] The bound must be enforced by a runtime watchdog timer that resets the system if exceeded.
    \item[\textbf{B)}] The bound must be statically verifiable, meaning it can be proven without executing the code.
    \item[\textbf{C)}] The bound must be dynamically calculated based on the available stack memory at the start of the loop.
    \item[\textbf{D)}] The bound must be manually certified by two independent engineers during the code review process.
\end{enumerate}
\answerspace{10}
\answer {
    B
}

\vspace{5mm}
\question[3] Which of the following best explains the reasoning behind the LOC-3 compliance rule concerning the use of assertions in flight software, requiring the software as a whole and each module within it to reach a minimal assertion density of 2\%?
\begin{enumerate}
	\item[\textbf{A)}] The rule aims to make the code work.
	\item[\textbf{B)}] The rule ensures that the code contains enough comments and documentation.
	\item[\textbf{C)}] The rule helps to improve software reliability by catching potential errors early and reducing residual defect densities.
	\item[\textbf{D)}] The rule enforces a strict coding style that prioritizes performance over maintainability.
\end{enumerate}
\answerspace{10}
\answer {
	C
}

\vspace{5mm}
\question[3] LOC-4 represents the highest target level for mission-critical code and restricts powerful features like function pointers. Despite this high level of strictness, how does the text characterize the \textit{quantity} of rules required to reach this level?
\begin{enumerate}
    \item[\textbf{A)}] It requires complying with hundreds of specific compiler-extension constraints.
    \item[\textbf{B)}] The cumulative number of coding rules remains relatively low, with no more than 31 risk-related rules.
    \item[\textbf{C)}] It requires a significant increase in rules, specifically doubling the number of assertions required by LOC-3.
    \item[\textbf{D)}] It introduces a complex set of over 100 new syntax restrictions to prevent pointer indirection.
\end{enumerate}
\answerspace{10}
\answer {
    B
}


\if 0 %%%%%%%%%%%%%%%%%%%%%%%%%%%%%%%%%%%%%%%%%
	\question[10]
	\answerspace{30}
	\answer{}

	\question
	\begin{parts}
		\part[6]
		\answerspace{20}
		\answer{}

		\part[2]
		\answerspace{10}
		\answer{}

	\end{parts}
\fi %%%%%%%%%%%%%%%%%%%%%%%%%%%%%%%%%%%%%%%%%

	\newpage
	\begin{center}
		[EMPTY PAGE]
	\end{center}

\end{questions}

\end{document}
